% Gerado por regressao/06_tabelas_latex.R (projeto Modelagem). Não editar à mão.
\begin{landscape}
\begin{table}[htbp]
\caption{Robustez R9 — só municípios que informam a função em todos os anos}\label{tab:robustez_R9}
\centering
\footnotesize
\setlength{\tabcolsep}{3pt}
\begin{tabular}{@{}>{\raggedright\arraybackslash}p{0.25\linewidth}>{\centering\arraybackslash}p{72pt}>{\centering\arraybackslash}p{72pt}>{\centering\arraybackslash}p{72pt}@{}}
\hline
\textbf{Variável} & \textbf{\hspace{0pt}Transporte} & \textbf{\hspace{0pt}Agricultura} & \textbf{\hspace{0pt}Comunicações} \\ \hline
\multicolumn{4}{@{}l}{\textit{Modelo A (ciclo eleitoral)}} \\
Ano eleitoral & 48,686\textsuperscript{***} & 0,077 & $-$0,300\textsuperscript{*} \\
 & (12,646) & (2,190) & (0,179) \\
Ano pré-eleitoral & 8,687 & 5,424 & 1,145\textsuperscript{***} \\
 & (9,526) & (3,997) & (0,150) \\
\hline
Observações & 38.194 & 54.221 & 4.237 \\
Municípios & 2.999 & 4.264 & 330 \\
R\textsuperscript{2} & 0,217 & 0,190 & 0,044 \\[6pt]
\multicolumn{4}{@{}l}{\textit{Modelo B (coalizões)}} \\
Prefeito na coalizão & 2,854 & $-$0,711 & 0,150 \\
do governador & (3,499) & (1,179) & (0,497) \\
Prefeito na coalizão & $-$17,037\textsuperscript{***} & $-$5,234\textsuperscript{***} & 0,241 \\
do presidente & (4,104) & (1,565) & (0,692) \\
\hline
Observações & 38.194 & 54.221 & 4.237 \\
Municípios & 2.999 & 4.264 & 330 \\
R\textsuperscript{2} & 0,106 & 0,081 & 0,022 \\
\hline
\end{tabular}
\fonte{Elaboração própria.}
\nota{Modelo A: efeito fixo de município, erros-padrão de Driscoll-Kraay. Modelo B: efeitos fixos de município e de ano, erros-padrão agrupados por município. Erros-padrão entre parênteses. \textsuperscript{***}p$<$0,01; \textsuperscript{**}p$<$0,05; \textsuperscript{*}p$<$0,1.}
\end{table}
\end{landscape}
